\documentclass[10pt,a4paper]{amsart}
\usepackage[margin=19mm]{geometry}
\usepackage{amsmath,amssymb,mathtools}
\usepackage[scr=rsfs,cal=euler]{mathalfa}
\usepackage{xfrac}
\usepackage[unicode,hidelinks,bookmarks=false]{hyperref}
\newcommand{\ie}{i.e.\ }
\newcommand{\cf}{cf.\ }
\newcommand{\resp}{respectively\ }
\newcommand{\Subsection}{\subsection}
\providecommand{\nobreakdash}{\nobreak\hbox{-}}
\setlength{\emergencystretch}{2em}
\multlinegap0pt
\usepackage{mdframed}
\usepackage{mdframed}
\usepackage{mdframed}
\usepackage{mdframed}
\usepackage{amssymb}
\usepackage{bm}
\usepackage{float}
\usepackage{mdframed}
\newtheorem{theorem}{Theorem}
\newtheorem{lemma}{Lemma}
\title{The Cauchy problem for the improved Boussinesq equation with spatially quasi-periodic initial data}
\author{Zhiqiang Wan, Wenji Wu, Heng Zhang}
\date{Source: arXiv:2605.07669v1 (CC BY 4.0)}
\begin{document}
\maketitle

\noindent\emph{Source and licence.} The Cauchy problem for the improved Boussinesq equation with spatially quasi-periodic initial data. Version arXiv:2605.07669v1 (CC BY 4.0), distributed by arXiv under the Creative Commons Attribution 4.0 International licence, http://creativecommons.org/licenses/by/4.0/, as recorded in the arXiv licence metadata for this exact version and shown on the abstract page https://arxiv.org/abs/2605.07669v1. Full licence notice: \texttt{licenses/arxiv-2605.07669-CC-BY-4.0.txt}.\par\smallskip

\noindent\textit{Editorial scope.} Counted unit: the article's theorem decay (source0.tex lines 198--232). The article's complete original proof of it is delivered with it (source0.tex lines 1664--2104, one \texttt{proof} environment of 441 lines, which the article places away from the statement). No mathematics is written, repaired or re-derived: the statement and the proof are the article's own lines, in the article's own notation, and no retained line is substituted or re-flowed.  Retained uncounted as the source-local context the proof needs: lines 149--152; lines 154--158; lines 451--451; lines 1082--1087; lines 1299--1314; lines 1301--1307; lines 1435--1457.  Nothing was omitted from any retained span, so the package carries no omission record.  No retained line contains a citation, so the wrapper carries no bibliography and no bibliography entry.  No retained line contains a figure environment, an image-inclusion command or drawing code, so no image is delivered and the package declares no asset dependency.  Editorial formatting only, in addition: an AMS \texttt{amsart} preamble that restates the article's own declarations and macros used by the retained lines, namely the environments \texttt{theorem}, \texttt{lemma} on separate counters, together with the standard packages those lines need.  The retained environments print in this document as Theorem 1, Lemma 1 in order of appearance, whereas the article prints the same items with the numbering of its own full document.  The complete native source of the article as distributed in the arXiv e-print for this exact version is bundled unchanged as \texttt{sources/200105/source0.tex}.
\begin{equation}\label{eq:IBq-intro}
    u_{tt}-u_{xx}-u_{xxtt}-(u^2)_{xx}=0,
    \qquad (t,x)\in \mathbb{R}\times\mathbb{R},
\end{equation}

\begin{equation}\label{eq:IBq-data-intro}
    u(0,x)=\sum_{n\in\mathbb{Z}^{\nu}}c(n)e^{i\langle n,\omega\rangle x},
    \qquad
    u_t(0,x)=\sum_{n\in\mathbb{Z}^{\nu}}d(n)e^{i\langle n,\omega\rangle x},
\end{equation}

\subsection{Reduction}\label{sec:reduction}

\begin{equation}\label{eq:exp-decay-initial}
|c(n)|\leq A e^{-\rho |n|},
\qquad
|d(n)|\leq A e^{-\rho |n|},
\qquad n\in \mathbb{Z}^{\nu},
\end{equation}

\begin{lemma}\label{lem:uniform-exp-decay}
Assume \eqref{eq:exp-decay-initial}. Set
\begin{equation}\label{eq:M-B-L}
M:=\max\{1,Ab_{\rho}\},
\qquad
B:=2M,
\qquad
L:=\frac{1}{5M}.
\end{equation}
Then, for every $n\in \mathbb{Z}^{\nu}$,
\begin{equation}\label{eq:uniform-exp-decay}
\sup_{0\leq t\leq L}\ \sup_{k\geq 0}\ |c_k(t,n)|
\leq
B e^{-\frac{\rho}{2}|n|}.
\end{equation}
\end{lemma}

\begin{equation}
M:=\max\{1,Ab_{\rho}\},
\qquad
B:=2M,
\qquad
L:=\frac{1}{5M}.
\end{equation}

\begin{lemma}\label{lem:cauchy-estimate}
Let $B$ and $L$ be as in \eqref{eq:M-B-L}. Then, for every $k\geq 1$, every
$0\leq t\leq L$, and every $n\in \mathbb{Z}^{\nu}$, one has
\begin{equation}\label{eq:cauchy-estimate-1}
|c_k(t,n)-c_{k-1}(t,n)|
\leq
\frac{2^{k-1}B^{k+1}t^k}{k!}\,\mathcal{E}_{k+1}(n).
\end{equation}
Moreover, setting
\begin{equation*}
\widetilde b_{\rho}:=(12\rho^{-1})^{\nu},
\end{equation*}
we also have
\begin{equation}\label{eq:cauchy-estimate-2}
|c_k(t,n)-c_{k-1}(t,n)|
\leq
\frac{B\widetilde b_{\rho}}{2}\,
\frac{(2B\widetilde b_{\rho}t)^k}{k!}\,
e^{-\frac{\rho}{4}|n|}.
\end{equation}
In particular, $\{c_k(t,n)\}_{k\geq 0}$ is a Cauchy sequence on
$[0,L]\times \mathbb{Z}^{\nu}$.
\end{lemma}

\begin{theorem}\label{thm:exponential decay}
Let $\omega\in \mathbb{R}^{\nu}$ be non-resonant, and consider the quasi-periodic
Cauchy problem for \eqref{eq:IBq-intro}
with initial data \eqref{eq:IBq-data-intro}.
Assume that there exist $A>0$ and $0<\rho\leq 1$ such that
\[
|c(n)|\leq A e^{-\rho |n|},
\qquad
|d(n)|\leq A e^{-\rho |n|},
\qquad n\in \mathbb{Z}^{\nu}.
\]
Set
\[
b_{\rho}:=(6\rho^{-1})^{\nu},
\qquad
M:=\max\{1,Ab_{\rho}\},
\qquad
B:=2M,
\qquad
L:=\frac{1}{5M}.
\]
Then, on the time interval $[0,L]$, the quasi-periodic Cauchy problem admits a classical spatially
quasi-periodic solution of the form
\[
u(t,x)=\sum_{n\in \mathbb{Z}^{\nu}} \widehat{u}(t,n)e^{i\langle n,\omega\rangle x}
\]
such that
\[
|\widehat{u}(t,n)|\leq B e^{-\frac{\rho}{2}|n|},
\qquad
0\leq t\leq L,\ \ n\in \mathbb{Z}^{\nu}.
\]
Moreover, this solution is unique among all spatially quasi-periodic solutions
satisfying the above exponential decay estimate.
\end{theorem}

\begin{proof}[Proof of Theorem \ref{thm:exponential decay}]
We divide the proof into three steps.

\medskip
\noindent
\textbf{Step 1: Construction of the limit Fourier coefficients.}
By Lemma~\ref{lem:cauchy-estimate}, the Picard sequence
$\{c_k(t,n)\}_{k\geq 0}$ is a Cauchy sequence on $[0,L]\times \mathbb{Z}^{\nu}$.
Hence there exists a limit function, denoted by $c^{\dagger}(t,n)$, such that
\[
c_k(t,n)\longrightarrow c^{\dagger}(t,n)
\qquad \text{uniformly on } [0,L]\times \mathbb{Z}^{\nu}.
\]
Moreover, by Lemma~\ref{lem:uniform-exp-decay},
\[
|c_k(t,n)|\leq B e^{-\frac{\rho}{2}|n|}
\qquad
\text{for all } k\geq 0,\ 0\leq t\leq L,\ n\in \mathbb{Z}^{\nu}.
\]
Passing to the limit yields
\begin{equation}\label{eq:limit-decay}
|c^{\dagger}(t,n)|\leq B e^{-\frac{\rho}{2}|n|},
\qquad
0\leq t\leq L,\ \ n\in \mathbb{Z}^{\nu}.
\end{equation}

We claim that $c^{\dagger}$ satisfies the integral equation
\begin{equation}\label{eq:limit-integral-equation}
c^{\dagger}(t,n)
=
c_0(t,n)
-
\int_0^t \Phi_n(t-\tau)
\sum_{\substack{n_1,n_2\in \mathbb{Z}^{\nu}\\ n_1+n_2=n}}
c^{\dagger}(\tau,n_1)c^{\dagger}(\tau,n_2)\,d\tau,
\end{equation}
where
\[
c_0(t,n)=G_n(t)c(n)+K_n(t)d(n).
\]
Indeed, for every fixed $(t,n)\in [0,L]\times \mathbb{Z}^{\nu}$,
\[
c_k(t,n)
=
c_0(t,n)
-
\int_0^t \Phi_n(t-\tau)
\sum_{\substack{n_1,n_2\in \mathbb{Z}^{\nu}\\ n_1+n_2=n}}
c_{k-1}(\tau,n_1)c_{k-1}(\tau,n_2)\,d\tau.
\]
Since $c_{k-1}\to c^{\dagger}$ uniformly and
\[
|c_{k-1}(\tau,n_1)c_{k-1}(\tau,n_2)|
\leq
B^2 e^{-\frac{\rho}{2}|n_1|}e^{-\frac{\rho}{2}|n_2|},
\]
the series over $n_1+n_2=n$ is absolutely convergent and dominated by a summable
majorant independent of $k$ and $\tau$. Therefore, by dominated convergence, we may
pass to the limit in the Picard iteration and obtain
\eqref{eq:limit-integral-equation}.

\medskip
\noindent
\textbf{Step 2: Existence of a classical solution.}
Define
\begin{equation*}
u^{\dagger}(t,x)
:=
\sum_{n\in \mathbb{Z}^{\nu}} c^{\dagger}(t,n)e^{i\langle n,\omega\rangle x}.
\end{equation*}
By \eqref{eq:limit-decay}, the series for $u^{\dagger}$ converges absolutely and
uniformly on $[0,L]\times \mathbb{R}$.

For every fixed $n\in \mathbb{Z}^{\nu}$, the function $t\mapsto c_k(t,n)$ is continuous
for every $k\geq 0$ by construction of the Picard sequence. Since
$c_k(\cdot,n)\to c^{\dagger}(\cdot,n)$ uniformly on $[0,L]$, it follows that
\begin{equation}\label{eq:c-dagger-continuous}
c^{\dagger}(\cdot,n)\in C([0,L]),
\qquad \forall\, n\in \mathbb{Z}^{\nu}.
\end{equation}

Next define, for $0\leq t\leq L$ and $n\in \mathbb{Z}^{\nu}$,
\begin{equation}\label{eq:nonlinear-coefficients}
\mathcal{N}^{\dagger}(t,n)
:=
\sum_{\substack{n_1,n_2\in \mathbb{Z}^{\nu}\\ n_1+n_2=n}}
c^{\dagger}(t,n_1)c^{\dagger}(t,n_2).
\end{equation}
By \eqref{eq:limit-decay},
\[
|\mathcal{N}^{\dagger}(t,n)|
\leq
B^2
\sum_{\substack{n_1,n_2\in \mathbb{Z}^{\nu}\\ n_1+n_2=n}}
e^{-\frac{\rho}{2}|n_1|}e^{-\frac{\rho}{2}|n_2|}.
\]
Arguing as in Lemma~\ref{lem:cauchy-estimate}, we obtain
\begin{equation}\label{eq:N-dagger-decay}
|\mathcal{N}^{\dagger}(t,n)|
\leq
B^2 \widetilde b_{\rho}^{\,2} e^{-\frac{\rho}{4}|n|},
\qquad
0\leq t\leq L,\ \ n\in \mathbb{Z}^{\nu},
\end{equation}
where $\widetilde b_{\rho}=(12\rho^{-1})^{\nu}$.

For each fixed $n$, the series in \eqref{eq:nonlinear-coefficients} is uniformly
absolutely convergent on $[0,L]$ by \eqref{eq:limit-decay}. Since each summand
$t\mapsto c^{\dagger}(t,n_1)c^{\dagger}(t,n_2)$ is continuous by
\eqref{eq:c-dagger-continuous}, it follows from the Weierstrass M-test that
\begin{equation*}
\mathcal{N}^{\dagger}(\cdot,n)\in C([0,L]),
\qquad \forall\, n\in \mathbb{Z}^{\nu}.
\end{equation*}

Now \eqref{eq:limit-integral-equation} is precisely the Duhamel formula for the linear
equation
\begin{equation}\label{eq:coefficient-ode-final}
\partial_t^2 c^{\dagger}(t,n)
+\Omega(n)^2 c^{\dagger}(t,n)
=
-\Omega(n)^2 \mathcal{N}^{\dagger}(t,n),
\end{equation}
with initial conditions
\begin{equation}\label{eq:coefficient-initial-final}
c^{\dagger}(0,n)=c(n),
\qquad
\partial_t c^{\dagger}(0,n)=d(n).
\end{equation}
Since $\mathcal{N}^{\dagger}(\cdot,n)\in C([0,L])$, standard ODE theory for the forced
harmonic oscillator yields
\begin{equation*}
c^{\dagger}(\cdot,n)\in C^2([0,L]),
\qquad \forall\, n\in \mathbb{Z}^{\nu},
\end{equation*}
and \eqref{eq:coefficient-ode-final}--\eqref{eq:coefficient-initial-final} hold pointwise.

We next derive bounds for the time derivatives. From
\eqref{eq:coefficient-ode-final} and \eqref{eq:N-dagger-decay}, using
$0\leq \Omega(n)^2\leq 1$, we obtain
\begin{equation}\label{eq:ctt-bound}
|\partial_t^2 c^{\dagger}(t,n)|
\leq
|c^{\dagger}(t,n)|+|\mathcal{N}^{\dagger}(t,n)|
\leq
B e^{-\frac{\rho}{2}|n|}
+
B^2 \widetilde b_{\rho}^{\,2} e^{-\frac{\rho}{4}|n|}
\leq
C_1 e^{-\frac{\rho}{4}|n|}
\end{equation}
for some constant $C_1>0$.

Differentiating \eqref{eq:limit-integral-equation} once and using
$\partial_t \Phi_n(t)=\beta(n)\partial_t K_n(t)$, we find
\begin{align*}
\partial_t c^{\dagger}(t,n)
&=
\partial_t G_n(t)\,c(n)
+
\partial_t K_n(t)\,d(n)
-
\int_0^t \partial_t \Phi_n(t-\tau)\,\mathcal{N}^{\dagger}(\tau,n)\,d\tau .
\end{align*}
Since $|\partial_t G_n(t)|\leq 1$, $|\partial_t K_n(t)|\leq 1$, and
$|\partial_t \Phi_n(t)|\leq 1$, it follows from
\eqref{eq:N-dagger-decay} that
\begin{equation}\label{eq:ct-bound}
|\partial_t c^{\dagger}(t,n)|
\leq
A e^{-\rho|n|}
+
A e^{-\rho|n|}
+
L B^2 \widetilde b_{\rho}^{\,2} e^{-\frac{\rho}{4}|n|}
\leq
C_2 e^{-\frac{\rho}{4}|n|}
\end{equation}
for some constant $C_2>0$.

Let
\[
C_{\omega}:=\max_{1\leq j\leq \nu} |\omega_j|.
\]
Since $|n|=\sum_{j=1}^{\nu}|n_j|$, we have
\begin{equation*}
|\langle n,\omega\rangle|
=
\left|\sum_{j=1}^{\nu} n_j\omega_j\right|
\leq
\sum_{j=1}^{\nu} |n_j||\omega_j|
\leq
C_{\omega}|n|.
\end{equation*}
Therefore,
\[
1+|\langle n,\omega\rangle|^2 \leq (1+C_{\omega}^2)(1+|n|^2).
\]
Combining this with \eqref{eq:limit-decay}, \eqref{eq:N-dagger-decay},
\eqref{eq:ct-bound}, and \eqref{eq:ctt-bound}, we get
\[
\sum_{n\in \mathbb{Z}^{\nu}}
\big(1+|\langle n,\omega\rangle|^2\big)
\Big(
|c^{\dagger}(t,n)|+|\partial_t c^{\dagger}(t,n)|
+|\partial_t^2 c^{\dagger}(t,n)|+|\mathcal{N}^{\dagger}(t,n)|
\Big)
<\infty
\]
uniformly for $t\in [0,L]$. Hence the Fourier series defining
$u^{\dagger}$, $u_t^{\dagger}$, $u_{tt}^{\dagger}$, $u_{xx}^{\dagger}$,
$u_{xxtt}^{\dagger}$, and $(u^{\dagger})^2_{xx}$ all converge absolutely and uniformly
on $[0,L]\times \mathbb{R}$. Therefore termwise differentiation is justified, and
\eqref{eq:coefficient-ode-final} implies that $u^{\dagger}$ is a classical solution of
\eqref{eq:IBq-intro}. The initial conditions follow from \eqref{eq:coefficient-initial-final}.
Finally, by \eqref{eq:limit-decay}, the solution $u^{\dagger}$ satisfies
\eqref{eq:IBq-intro} with $\widehat u(t,n)=c^{\dagger}(t,n)$.

\medskip
\noindent
\textbf{Step 3: Uniqueness.}
Let
\[
v(t,x)=\sum_{n\in \mathbb{Z}^{\nu}} \widehat v(t,n)e^{i\langle n,\omega\rangle x},
\qquad
w(t,x)=\sum_{n\in \mathbb{Z}^{\nu}} \widehat w(t,n)e^{i\langle n,\omega\rangle x},
\]
be two spatially quasi-periodic classical solutions satisfying
\begin{equation}\label{eq:uniqueness-decay}
|\widehat v(t,n)|\leq B e^{-\frac{\rho}{2}|n|},
\qquad
|\widehat w(t,n)|\leq B e^{-\frac{\rho}{2}|n|},
\qquad
0\leq t\leq L,\ \ n\in \mathbb{Z}^{\nu}.
\end{equation}
By the same reduction argument as in Section~\ref{sec:reduction}, which is justified
here by the exponential decay \eqref{eq:uniqueness-decay}, both $\widehat v$ and
$\widehat w$ satisfy the coefficient equation
\[
\partial_t^2 a(t,n)+\Omega(n)^2 a(t,n)
=
-\Omega(n)^2
\sum_{\substack{n_1,n_2\in \mathbb{Z}^{\nu}\\ n_1+n_2=n}}
a(t,n_1)a(t,n_2),
\]
together with the same initial conditions. Equivalently, both satisfy the same integral
equation
\[
a(t,n)
=
G_n(t)a(0,n)+K_n(t)\partial_t a(0,n)
-
\int_0^t \Phi_n(t-\tau)
\sum_{\substack{n_1,n_2\in \mathbb{Z}^{\nu}\\ n_1+n_2=n}}
a(\tau,n_1)a(\tau,n_2)\,d\tau .
\]
Since $v$ and $w$ have the same initial data, we obtain
\begin{align}
\widehat v(t,n)-\widehat w(t,n)
&=
-\int_0^t \Phi_n(t-\tau)
\sum_{\substack{n_1,n_2\in \mathbb{Z}^{\nu}\\ n_1+n_2=n}}
\Big(
\widehat v(\tau,n_1)\widehat v(\tau,n_2)
-
\widehat w(\tau,n_1)\widehat w(\tau,n_2)
\Big)\,d\tau .
\label{eq:difference-integral}
\end{align}

Let
\[
\delta(t,n):=\widehat v(t,n)-\widehat w(t,n),
\qquad
0\leq t\leq L,\ \ n\in \mathbb{Z}^{\nu}.
\]
Then \eqref{eq:difference-integral} becomes
\begin{equation}\label{eq:difference-integral-delta}
\delta(t,n)
=
-\int_0^t \Phi_n(t-\tau)
\sum_{\substack{n_1,n_2\in \mathbb{Z}^{\nu}\\ n_1+n_2=n}}
\Big(
\delta(\tau,n_1)\widehat w(\tau,n_2)
+
\widehat v(\tau,n_1)\delta(\tau,n_2)
\Big)\,d\tau .
\end{equation}
Since \eqref{eq:uniqueness-decay} holds, we also have the crude bound
\begin{equation}\label{eq:delta-crude}
|\delta(t,n)|
\leq
|\widehat v(t,n)|+|\widehat w(t,n)|
\leq
2B e^{-\frac{\rho}{2}|n|},
\qquad
0\leq t\leq L,\ \ n\in \mathbb{Z}^{\nu}.
\end{equation}

We claim that, for every $k\geq 1$,
\begin{equation}\label{eq:uniqueness-claim}
|\delta(t,n)|
\leq
\frac{2^{k+1}B^{k+1}t^k}{k!}\,\mathcal{E}_{k+1}(n),
\qquad
0\leq t\leq L,\ \ n\in \mathbb{Z}^{\nu},
\end{equation}
where
\[
\mathcal{E}_{m}(n)
:=
\sum_{\substack{n_1,\dots,n_m\in \mathbb{Z}^{\nu}\\ n_1+\cdots+n_m=n}}
\prod_{j=1}^{m} e^{-\frac{\rho}{2}|n_j|}.
\]

For $k=1$, by \eqref{eq:difference-integral-delta}, \eqref{eq:uniqueness-decay},
and \eqref{eq:delta-crude}, we obtain
\begin{align*}
|\delta(t,n)|
&\leq
\int_0^t |\Phi_n(t-\tau)|
\sum_{\substack{n_1,n_2\in \mathbb{Z}^{\nu}\\ n_1+n_2=n}}
\Big(
|\delta(\tau,n_1)|\,|\widehat w(\tau,n_2)|
+
|\widehat v(\tau,n_1)|\,|\delta(\tau,n_2)|
\Big)\,d\tau \\
&\leq
\int_0^t
\sum_{\substack{n_1,n_2\in \mathbb{Z}^{\nu}\\ n_1+n_2=n}}
\Big(
2B e^{-\frac{\rho}{2}|n_1|}\cdot B e^{-\frac{\rho}{2}|n_2|}
+
B e^{-\frac{\rho}{2}|n_1|}\cdot 2B e^{-\frac{\rho}{2}|n_2|}
\Big)\,d\tau \\
&=
4B^2 t
\sum_{\substack{n_1,n_2\in \mathbb{Z}^{\nu}\\ n_1+n_2=n}}
e^{-\frac{\rho}{2}|n_1|}e^{-\frac{\rho}{2}|n_2|} \\
&=
\frac{2^{2}B^{2}t}{1!}\,\mathcal{E}_{2}(n).
\end{align*}
Hence \eqref{eq:uniqueness-claim} holds for $k=1$.

Now let $k\geq 2$ and assume that \eqref{eq:uniqueness-claim} holds with $k-1$
in place of $k$, that is,
\[
|\delta(\tau,m)|
\leq
\frac{2^{k}B^{k}\tau^{k-1}}{(k-1)!}\,\mathcal{E}_{k}(m),
\qquad
0\leq \tau\leq L,\ \ m\in \mathbb{Z}^{\nu}.
\]
Using \eqref{eq:difference-integral-delta}, \eqref{eq:uniqueness-decay}, and the
induction hypothesis, we get
\begin{align*}
|\delta(t,n)|
&\leq
\int_0^t
\sum_{\substack{n_1,n_2\in \mathbb{Z}^{\nu}\\ n_1+n_2=n}}
|\delta(\tau,n_1)|\,|\widehat w(\tau,n_2)|\,d\tau \\
&\quad+
\int_0^t
\sum_{\substack{n_1,n_2\in \mathbb{Z}^{\nu}\\ n_1+n_2=n}}
|\widehat v(\tau,n_1)|\,|\delta(\tau,n_2)|\,d\tau \\
&\leq
\int_0^t
\sum_{\substack{n_1,n_2\in \mathbb{Z}^{\nu}\\ n_1+n_2=n}}
\frac{2^{k}B^{k}\tau^{k-1}}{(k-1)!}\,\mathcal{E}_{k}(n_1)\,
B e^{-\frac{\rho}{2}|n_2|}\,d\tau \\
&\quad+
\int_0^t
\sum_{\substack{n_1,n_2\in \mathbb{Z}^{\nu}\\ n_1+n_2=n}}
B e^{-\frac{\rho}{2}|n_1|}\,
\frac{2^{k}B^{k}\tau^{k-1}}{(k-1)!}\,\mathcal{E}_{k}(n_2)\,d\tau .
\end{align*}
By the definition of $\mathcal{E}_{k}$ and a relabeling of the frequency variables,
\[
\sum_{\substack{n_1,n_2\in \mathbb{Z}^{\nu}\\ n_1+n_2=n}}
\mathcal{E}_{k}(n_1)e^{-\frac{\rho}{2}|n_2|}
=
\mathcal{E}_{k+1}(n),
\]
and similarly
\[
\sum_{\substack{n_1,n_2\in \mathbb{Z}^{\nu}\\ n_1+n_2=n}}
e^{-\frac{\rho}{2}|n_1|}\mathcal{E}_{k}(n_2)
=
\mathcal{E}_{k+1}(n).
\]
Therefore
\begin{align*}
|\delta(t,n)|
&\leq
\frac{2^{k}B^{k+1}}{(k-1)!}
\int_0^t \tau^{k-1}\,d\tau \, \mathcal{E}_{k+1}(n)
+
\frac{2^{k}B^{k+1}}{(k-1)!}
\int_0^t \tau^{k-1}\,d\tau \, \mathcal{E}_{k+1}(n) \\
&=
\frac{2^{k}B^{k+1}t^k}{k!}\,\mathcal{E}_{k+1}(n)
+
\frac{2^{k}B^{k+1}t^k}{k!}\,\mathcal{E}_{k+1}(n) \\
&=
\frac{2^{k+1}B^{k+1}t^k}{k!}\,\mathcal{E}_{k+1}(n).
\end{align*}
Thus \eqref{eq:uniqueness-claim} holds for all $k\geq 1$.

By the same argument as in Lemma~\ref{lem:cauchy-estimate},
\[
\mathcal{E}_{k+1}(n)
\leq
\widetilde b_{\rho}^{\,k+1} e^{-\frac{\rho}{4}|n|}.
\]
Hence \eqref{eq:uniqueness-claim} implies
\[
|\delta(t,n)|
\leq
2B\widetilde b_{\rho}\,
\frac{(2B\widetilde b_{\rho}t)^k}{k!}
e^{-\frac{\rho}{4}|n|},
\qquad
k\geq 1.
\]
For every fixed $(t,n)\in [0,L]\times \mathbb{Z}^{\nu}$, the right-hand side tends to
zero as $k\to\infty$. Therefore
\[
\delta(t,n)\equiv 0,
\qquad
0\leq t\leq L,\ \ n\in \mathbb{Z}^{\nu}.
\]
That is,
\[
\widehat v(t,n)\equiv \widehat w(t,n),
\qquad
0\leq t\leq L,\ \ n\in \mathbb{Z}^{\nu}.
\]
Hence $v\equiv w$, proving uniqueness.

Combining the three steps, the proof of Theorem~\ref{thm:exponential decay} is complete.
\end{proof}

\end{document}
